
\documentclass[11pt,a4paper]{article}

% ---- geometry ----
\usepackage[margin=2.4cm]{geometry}

% ---- encoding and fonts ----
\usepackage[utf8]{inputenc}
\usepackage[T1]{fontenc}
\usepackage{lmodern}
\usepackage{microtype}

% ---- maths ----
\usepackage{amsmath,amssymb,amsthm,mathtools}
\usepackage{bm}

% ---- figures and tables ----
\usepackage{graphicx}
\graphicspath{{figures/}}
\usepackage{booktabs}
\usepackage{array}
\usepackage{float}
\usepackage[section]{placeins}

% ---- float placement ----
\renewcommand{\topfraction}{0.9}
\renewcommand{\bottomfraction}{0.9}
\renewcommand{\textfraction}{0.1}
\renewcommand{\floatpagefraction}{0.7}
\setcounter{topnumber}{4}
\setcounter{bottomnumber}{4}
\setcounter{totalnumber}{8}

% ---- references and links ----
\usepackage[colorlinks=true,linkcolor=blue!60!black,
            citecolor=green!50!black,urlcolor=blue!70!black]{hyperref}
\usepackage{bookmark}
\usepackage[nameinlink,capitalize]{cleveref}

% ---- bibliography ----
\usepackage[numbers,sort&compress]{natbib}

% ---- misc ----
\usepackage{enumitem}
\usepackage{xcolor}

% ---- theorem-like environments ----
\usepackage{aliascnt}
\theoremstyle{definition}

\newtheorem{theorem}{Theorem}[section]

\newaliascnt{proposition}{theorem}
\newtheorem{proposition}[proposition]{Proposition}
\aliascntresetthe{proposition}

\newaliascnt{lemma}{theorem}
\newtheorem{lemma}[lemma]{Lemma}
\aliascntresetthe{lemma}

\newaliascnt{corollary}{theorem}
\newtheorem{corollary}[corollary]{Corollary}
\aliascntresetthe{corollary}

\newaliascnt{definition}{theorem}
\newtheorem{definition}[definition]{Definition}
\aliascntresetthe{definition}

\newaliascnt{remark}{theorem}
\newtheorem{remark}[remark]{Remark}
\aliascntresetthe{remark}

% ---- cleveref names ----
\crefname{proposition}{Proposition}{Propositions}
\Crefname{proposition}{Proposition}{Propositions}
\crefname{lemma}{Lemma}{Lemmas}
\Crefname{lemma}{Lemma}{Lemmas}
\crefname{corollary}{Corollary}{Corollaries}
\Crefname{corollary}{Corollary}{Corollaries}
\crefname{definition}{Definition}{Definitions}
\Crefname{definition}{Definition}{Definitions}
\crefname{remark}{Remark}{Remarks}
\Crefname{remark}{Remark}{Remarks}

% ---- custom commands ----
\newcommand{\khat}{\hat{\bm{k}}}
\newcommand{\nhat}{\hat{\bm{n}}}
\newcommand{\rr}{\mathbf{r}}
\newcommand{\EE}{\mathbf{E}}
\newcommand{\Sinc}{S_{\mathrm{inc}}}
\newcommand{\Sab}{S_{\mathrm{ab}}}
\newcommand{\Teff}{T_{\mathrm{eff}}}
\newcommand{\Tavg}{T_{\mathrm{avg}}}
\newcommand{\Aab}{A_{\mathrm{ab}}}
\newcommand{\Aperp}{A_{\perp}}
\newcommand{\SARwb}{\mathrm{SAR}_{\mathrm{wb}}}
\newcommand{\ntilde}{\tilde{n}}
\DeclareMathOperator{\ReLU}{ReLU}
\DeclareMathOperator{\RE}{Re}
\newcommand{\Tbar}{\bar{T}}
\newcommand{\diff}{\mathrm{d}}

% ---- Part headings ----
\newcommand{\partheading}[2]{%
  \clearpage
  \phantomsection
  \addcontentsline{toc}{part}{#1}%
  \begin{center}
    \vspace*{1.5cm}
    {\LARGE\bfseries #1}\\[0.6cm]
    {\large #2}\\[1.5cm]
  \end{center}
}

% ---- title block ----
\title{\textbf{Geometric Dosimetry:\\
Closed-Form Absorption Laws from 100\,MHz to 100\,GHz}}
\author{R.\ Wydaeghe}
\date{\today}

% ============================================================
\begin{document}
\maketitle

% ============================================================
\begin{abstract}
Brute-force electromagnetic simulation of absorbed power density on
human bodies at millimetre-wave frequencies ($>6$\,GHz) is intractable:
the skin depth (${\sim}\,0.5$\,mm) is five orders of magnitude smaller
than the body, requiring ${\sim}\,10^{12}$ mesh cells.
We show that for biological tissue the absorbed power density reduces
to a closed-form expression separating material response from geometry.
A pseudo-Brewster compensation between TE and TM polarisations
renders the angle-averaged Fresnel transmission nearly constant,
collapsing the material factor to a single scalar~$T_0$.  The remaining
geometric factor is the positive-part cosine of the surface normal, governed
by projected area, Cauchy's formula, and an exposure fraction identical
to ambient occlusion in computer graphics.  For multiple sources, the result
is a single-hidden-layer ReLU neural network, making the dosimetry chain
differentiable and compatible with ray-tracing engines.

Two generalisations extend the framework.  An absorption Stokes vector
handles arbitrary polarisation exactly; the worst-case polarisation
correction on a human body is 16\,\%, suppressed below 2.5\,\% in
multipath environments.  Replacing $T_0$ with an exact flux-averaged
value~$\Tbar$ makes total-power results exact at any frequency above
100\,MHz.

Part~III extends to coherent MIMO beamforming, where two well-bounded
approximations collapse the coherent depth integral to a squared norm
$\Sab(\mathbf{r})=\|\tilde{\mathbf{G}}(\mathbf{r})\mathbf{x}\|^2$.  An
exposure operator~$\mathbf{Q}$ encodes the body's complete absorption
response to any precoder and yields a closed-form exposure-constrained
beamformer.
\end{abstract}

\tableofcontents
\newpage

% ============================================================
\section*{Preface}
\addcontentsline{toc}{section}{Preface}
% ============================================================

I set out to interpret some FDTD results for millimetre-wave absorption
on the human body, a problem that by brute force requires trillions of
mesh cells.  To my surprise, the TE and TM Fresnel transmissions through
tissue compensate each other, causing the material factor to collapse to
a single constant.  This pseudo-Brewster compensation is a consequence
of the relatively high refractive index of living tissue.

Once the material factor becomes a constant, dosimetry becomes geometry.
The direction-averaged absorption satisfies Cauchy's projected-area
formula, first proved in 1841.  The self-shadowing correction is ambient
occlusion in computer graphics.  The multi-source formula is a ReLU
neural network whose weights come from a ray tracer.  At every turn, the
answer was a known result in another field.

\newpage

% ============================================================
\section{Introduction}\label{sec:introduction}
% ============================================================

\subsection{Problem and approach}

Millimetre-wave (mmWave) wireless systems operating above 6\,GHz
require regulatory evaluation of the absorbed power density~$\Sab$
on the human body~\cite{ICNIRP2020}.  At mmWave frequencies the skin
depth $\delta\sim 0.3$ to $0.5$\,mm is far smaller than any body
dimension, so all power that penetrates the surface is absorbed within a
thin surface layer.  Full-wave methods (FDTD, FEM) require a mesh
of order $10^{12}$~cells, far beyond present computational resources.

We exploit two physical simplifications specific to biological tissue.
First, the pseudo-Brewster compensation between TE and TM Fresnel
transmissions makes the angle-averaged transmission nearly independent
of incidence angle over the mmWave band.  Second, the small skin depth
confines all absorption to the surface, so only the surface geometry
matters.  Together they reduce the problem to a single formula:
\begin{equation}\label{eq:geometric-law}
  \boxed{\Sab(\rr)
    \;\approx\; \Sinc\;\cdot T_0\cdot
    \ReLU\!\bigl[\nhat(\rr)\cdot(-\khat)\bigr],}
\end{equation}
where $\Sinc$ is the incident power density, $T_0$ is the
\emph{normal-incidence} Fresnel power-absorption coefficient (a scalar
that depends only on tissue type and frequency), $\nhat(\rr)$ is the
outward surface normal at point $\rr$, $\khat$ is the direction of
propagation, and $\ReLU(x)=\max(0,x)$ gates the illuminated hemisphere.
All spatial variation is determined by the body's shape alone.

The result connects three fields.  In integral geometry, total absorbed
power equals $\Sinc\,T_0$ times the projected shadow area, and
direction-averaged absorption satisfies Cauchy's formula~\cite{Cauchy1841}.
In computer graphics, the self-shadowing factor is ambient
occlusion~\cite{Zhukov1998}.  For multiple incident waves, $\Sab$ is a
single-hidden-layer ReLU neural network~\cite{NairHinton2010}, making
the dosimetry chain differentiable.

\subsection{Scope and accuracy}

The core results apply from 100\,MHz to 100\,GHz.  The geometric
approximation \cref{eq:geometric-law} is optimal in the 10--60\,GHz
band, where the Fresnel error is 2 to 6\,\%.  Total-power and compliance
results extend to sub-6\,GHz via the exact flux-averaged
transmission~$\Tbar$ (\cref{sec:sub6}).  Comparison with Mie theory for
lossy spheres (the worst-case geometry for the framework) shows errors
of roughly 10\,\% for total absorbed power, dominated by diffraction into
the geometric shadow.  The uncertainty in tissue dielectric properties
is itself 20\,\%, so the framework error is comparable to the parametric
uncertainty of any simulation.  \Cref{fig:mie} quantifies both regimes.

\begin{figure}[!htbp]
\centering
\includegraphics{figures/mie_validation_corrected.pdf}
\caption{Validation against Mie theory for a lossy sphere (skin at 28\,GHz).
  Top: framework error vs.~size parameter $x=\pi d/\lambda$ for a
  single illumination direction.  The error approaches the Fresnel limit
  ($R=0.988$, orange) for large spheres; small spheres are dominated
  by diffraction.  Annotated vertical lines mark body-part diameters.
  Middle: error vs.~frequency for fixed diameters.  The mmWave band
  (20--60\,GHz, shaded) is where torso-sized bodies fall in the
  geometric-optics regime and errors stay below 10\,\%.
  Bottom: the sphere ratio $R=T_0/\Tbar$, which crosses unity at 40\,GHz.}
\label{fig:mie}
\end{figure}

\subsection{Relationship to prior work}

Whole-body SAR has relied on FDTD of anatomical voxel
models~\cite{Dimbylow2002,Christ2006}; no closed-form surface-based
framework existed before this work.  Bamba~\textit{et
al.}~\cite{Bamba2014} empirically fitted an absorption efficiency
$\eta\approx 0.5$; the present framework identifies $\eta$ with $T_0$
and explains its near-constancy as a consequence of pseudo-Brewster
compensation.  Azzam~\cite{Azzam2015} showed that for $|\ntilde|>2.5$
the unpolarised reflectance is nearly angle-independent; biological
tissue falls in this regime.  Cauchy's formula~\cite{Cauchy1841} and
ambient occlusion~\cite{Zhukov1998} are classical; neither has been
applied to dosimetry before this work.

% ============================================================
\partheading{Part~I: A Millimetre-Wave Framework}
{Sections~\ref{sec:local-law}--\ref{sec:compliance}: a self-contained
dosimetry framework for the 10--60\,GHz band.}
% ============================================================

% ============================================================
\section{Local absorption law}\label{sec:local-law}
% ============================================================

The starting point is the exact, pointwise expression for absorbed
power density at a body-surface point, valid before any approximation.

\subsection{Exact local law}

Consider a surface point $\rr$ with outward normal $\nhat(\rr)$.  An
incident plane wave with Poynting vector $\Sinc\,\khat$ hits the surface
at incidence cosine $\mu(\rr)\equiv\nhat(\rr)\cdot(-\khat)$.  The wave
decomposes locally into a TE component (electric field perpendicular to
the incidence plane) and a TM component (magnetic field perpendicular).
Let $T_s(\theta)$ and $T_p(\theta)$ be the Fresnel power-absorption
coefficients for TE and TM respectively, with local energy fractions
$|e_s|^2$ and $|e_p|^2$ ($|e_s|^2+|e_p|^2=1$).  Define the
\emph{effective transmission} as the polarisation-weighted average:
\begin{equation}\label{eq:Teff}
  \Teff(\rr) \;\equiv\; |e_s(\rr)|^2\,T_s(\theta(\rr)) + |e_p(\rr)|^2\,T_p(\theta(\rr)).
\end{equation}
Energy conservation then gives the \emph{exact absorption law}:
\begin{equation}\label{eq:Sab-exact}
  \boxed{\Sab(\rr) = \Sinc\;\Teff(\rr)\;\ReLU\!\bigl(\mu(\rr)\bigr).}
\end{equation}
The ReLU gate is exact: back-facing surface points ($\mu<0$) are in
shadow and absorb nothing.

\subsection{Polarisation decomposition}\label{sec:decomp}

It is useful to split $\Teff$ into a polarisation-independent part and a
correction.  Define the \emph{angle-averaged transmission}
$\Tavg(\theta)=\frac{1}{2}[T_s(\theta)+T_p(\theta)]$, the
\emph{polarisation splitting} $\Delta T(\theta)=T_p(\theta)-T_s(\theta)\ge
0$, and the \emph{local TM excess}
$q(\rr)=|e_p(\rr)|^2-|e_s(\rr)|^2\in[-1,1]$ (positive when the local
polarisation is more TM than TE).  Then:
\begin{equation}\label{eq:Teff-decomp}
  \Teff(\rr) = \Tavg\!\bigl(\theta(\rr)\bigr)
    + \frac{q(\rr)}{2}\,\Delta T\!\bigl(\theta(\rr)\bigr).
\end{equation}
The second term is the polarisation correction.  It vanishes for
circularly polarised or truly unpolarised illumination (where $\langle
q\rangle=0$ by symmetry), and is treated exactly in Part~II.  For
Part~I, we set $q=0$ and work with $\Tavg(\theta)$ alone.

% ============================================================
\section{Pseudo-Brewster compensation}\label{sec:pseudo-brewster}
% ============================================================

With $q=0$, the exact law reduces to
$\Sab=\Sinc\,\Tavg(\theta)\,\ReLU(\mu)$.  The approximation
$\Tavg(\theta)\approx T_0$ (the normal-incidence value) is the
\emph{pseudo-Brewster compensation}.

For a lossy dielectric, $T_s(\theta)$ decreases monotonically with
incidence angle while $T_p(\theta)$ first rises (pseudo-Brewster peak
near $\theta_{pB}\approx\arctan|\ntilde|$) then falls.  Their average
$\Tavg(\theta)$ remains close to $T_0$ over a broad angular range.
Azzam~\cite{Azzam2015} showed analytically that the unpolarised
reflectance $1-\Tavg$ is nearly angle-independent for $|\ntilde|>2.5$.
Biological tissue at mmWave has $|\ntilde|\approx 3$--6, precisely in
this regime.

For skin at 28\,GHz ($\ntilde\approx 4.49-1.79i$, $T_0=0.536$),
$\Tavg(\theta)$ varies by at most 5.6\,\% over
$\theta\in[0^\circ,75^\circ]$.  The approximation $\Tavg\approx T_0$
is optimal near 40\,GHz, where it is exact to four significant figures,
and introduces a conservative error (underestimate) below 40\,GHz
and a 3.5\,\% overestimate above 80\,GHz.  Below 10\,GHz the
approximation should be replaced by the exact averaged
transmission~$\Tbar$ (\cref{sec:sub6}).

\begin{table}[!htbp]
\centering
\caption{Pseudo-Brewster compensation across representative tissue types
  at 28\,GHz.  $|\ntilde|$ is the modulus of the complex refractive index;
  $\Tavg$ refers to the variation of $\Tavg(\theta)/T_0$ over
  $[0^\circ,75^\circ]$.}\label{tab:materials}
\begin{tabular}{lccr}
\toprule
Tissue & $|\ntilde|$ & $T_0$ & $\Tavg/T_0$ variation \\
\midrule
Skin   & 4.84 & 0.54 & $\le 5.6\,\%$ \\
Muscle & 5.62 & 0.48 & $\le 4.8\,\%$ \\
Fat    & 2.05 & 0.77 & $\le 8.2\,\%$ \\
Water  & 6.62 & 0.45 & $\le 3.9\,\%$ \\
\bottomrule
\end{tabular}
\end{table}

With $\Tavg\approx T_0$ the geometric absorption law
\cref{eq:geometric-law} follows immediately.  \Cref{fig:fresnel} plots
both transmittance curves and their average for skin at 28\,GHz.

\begin{figure}[!htbp]
\centering
\includegraphics[width=\linewidth]{figures/fresnel_curves.pdf}
\caption{Fresnel power transmittance for skin at 28\,GHz ($\ntilde=4.49-1.79i$).
  Left: $T_s$ (TE, falling), $T_p$ (TM, rising to pseudo-Brewster peak near
  $75^\circ$), and their average $\Tavg$ (dashed), which stays within 5.6\,\%
  of the normal-incidence value $T_0=0.539$ (dotted) up to $75^\circ$.
  Right: the corresponding Fresnel amplitude $|t|$ and phase $\arg(t)$ used
  in the coherent framework of Part~III.}
\label{fig:fresnel}
\end{figure}

% ============================================================
\section{Geometric framework}\label{sec:geometry}
% ============================================================

\subsection{Total absorbed power and the projected area}

Integrating \cref{eq:geometric-law} over the illuminated surface,
\begin{equation}\label{eq:Pabs-Aperp}
  P_{\mathrm{abs}}(\khat)
  = \Sinc\,T_0\,\Aperp(\khat),
\end{equation}
where $\Aperp(\khat)\equiv\int_{\Sigma_+}\cos\theta\,\diff A$ is the
\emph{projected area}: the shadow cast by the body perpendicular
to~$\khat$.  This is a purely geometric quantity, independent of all material
properties.

\subsection{Exposure fraction and Cauchy's formula}

For a non-convex body (such as a standing human), some surface points
are in self-shadow: illuminated from one direction but hidden behind
another body part.  The self-shadow is captured by the occlusion
factor $O(\rr,\khat)\in\{0,1\}$, which is 1 if and only if $\rr$ is
directly visible from direction~$\khat$.

\begin{definition}[Exposure fraction]\label{def:eta}
The \emph{exposure fraction} $\eta(\rr)$ is the cosine-weighted fraction
of the upper hemisphere from which $\rr$ is unoccluded:
\begin{equation}\label{eq:eta-def}
  \eta(\rr) \;\equiv\; \frac{1}{\pi}\int_{S^2}
  \ReLU\!\bigl[\nhat(\rr)\cdot(-\khat)\bigr]\;
  O(\rr,\khat)\;\diff\Omega(\khat).
\end{equation}
For convex bodies $\eta\equiv 1$; for concave regions $\eta<1$.
\end{definition}

The exposure fraction is mathematically identical to \emph{ambient
occlusion} in computer graphics~\cite{Zhukov1998}, making the dosimetry
map available from heavily optimised GPU implementations.

\begin{definition}[Absorption area]\label{def:Aab}
$\Aab\equiv\int_\Sigma\eta(\rr)\,\diff A$: the surface area weighted by
exposure fraction.  For a convex body $\Aab=A$; for the Thelonious
human phantom (23\,826 triangles) $\Aab/A\approx 0.87$.
\end{definition}

\begin{theorem}[Generalised Cauchy formula]\label{thm:gen-cauchy}
The direction-averaged total absorbed power satisfies
\begin{equation}\label{eq:gen-cauchy}
  \boxed{\langle P_{\mathrm{abs}}\rangle
    = \Sinc\,T_0\,\frac{\Aab}{4}.}
\end{equation}
\end{theorem}

\begin{proof}
Integrate $P_{\mathrm{abs}}(\khat)=\Sinc\,T_0\,\Aperp(\khat)$ over
$\khat\in S^2$.  Exchange the order of integration (Fubini).  The
directional integral $\frac{1}{4\pi}\int_{S^2}[\nhat\cdot(-\khat)]_+\,O\,\diff\Omega
=\eta(\rr)/4$ by the definition of $\eta$.  Integrating over the
surface gives $\Aab/4$.
\end{proof}

The classical Cauchy formula ($\langle\Aperp\rangle=A/4$ for convex
bodies~\cite{Cauchy1841}) is the special case $\eta\equiv 1$.
\Cref{fig:eta3d,fig:sab3d} show the exposure fraction and the resulting
absorbed power density map on the Thelonious phantom.

\begin{figure}[!htbp]
\centering
\begin{minipage}[b]{0.47\linewidth}
  \centering
  \includegraphics[width=\linewidth]{figures/eta_3d_phantom.png}
  \caption{Exposure fraction $\eta(\mathbf{r})$ on the Thelonious phantom.
    Most of the surface is fully exposed ($\eta=1$, red); reduced values
    appear only in the armpits and inner thighs where self-shadowing is
    significant.  The area-weighted mean is $\Aab/A=0.87$.}
  \label{fig:eta3d}
\end{minipage}\hfill
\begin{minipage}[b]{0.47\linewidth}
  \centering
  \includegraphics[width=\linewidth]{figures/sab_3d_4cm2.png}
  \caption{Absorbed power density $\Sab(\mathbf{r})$ (W/m$^2$) averaged over 4 cm$^2$ predicted
    by the geometric absorption law for a single plane wave incident from
    the front.  The cosine-law pattern is clearly visible: the torso and
    head receive the highest flux; limbs and the back face zero.}
  \label{fig:sab3d}
\end{minipage}
\end{figure}

\subsection{Absorption directivity}

\begin{definition}[Absorption directivity]
$D(\khat)\equiv 4\Aperp(\khat)/\Aab$.  This is the ratio of the projected
area from direction~$\khat$ to the direction-averaged projected area.
For an isotropic radiator $D=1$; for the Thelonious phantom
$D_{\max}\approx 1.22$ (side-on) and $D_{\min}\approx 0.45$
(overhead).
\end{definition}

\begin{proposition}[Factorised compliance quantity]\label{prop:factorised}
\begin{equation}\label{eq:Pabs-D}
  \boxed{P_{\mathrm{abs}}(\khat) = \Sinc\,T_0\,\frac{\Aab}{4}\;D(\khat).}
\end{equation}
The electromagnetic environment enters only through $\Sinc$; the tissue
through~$T_0$; and the body geometry through $\Aab$ and $D(\khat)$,
precomputable from the body mesh.
\end{proposition}

Since $D(\khat)$ is smooth on $S^2$, it admits a low-order spherical
harmonic expansion: $L\le 6$ (49~coefficients) achieves an RMS error
below 1\,\% for typical human phantoms.

% ============================================================
\section{Compliance}\label{sec:compliance}
% ============================================================

The ICNIRP limit on whole-body SAR is 0.08\,W/kg for general
public~\cite{ICNIRP2020}.  Using $\SARwb=P_{\mathrm{abs}}/m$
together with \cref{eq:Pabs-D} gives the worst-case compliance
criterion (using $D\le 2$, the hemisphere bound):
\begin{equation}\label{eq:compliance-worst}
  \boxed{\Sinc < \frac{0.16\,m}{T_0\,\Aab}.}
\end{equation}
For a convex body one may further bound $\Aab\le A$, replacing
$\Aab$ by the total surface area.  The threshold depends only on
body mass $m$, absorption area~$\Aab$, and the tissue-frequency
property~$T_0$.  Via the Du~Bois formula~\cite{DuBois1916}, the
ratio $m/\Aab$ scales approximately as $\mathrm{BMI}^{0.575}h^{0.425}$,
so taller and heavier individuals tolerate more incident power.
Children, whose smaller $m/\Aab$ ratio lowers the threshold, are
the binding case for population-level reference levels.
\Cref{fig:error-budget} summarises the relative magnitude of each
error source across best- and worst-case scenarios.

\begin{figure}[!htbp]
\centering
\includegraphics[width=\linewidth]{figures/error_budget_comprehensive.pdf}
\caption{Error budget for the geometric absorption law.  (a)~Best case:
  28\,GHz, torso illumination, $N=20$ multipath paths.  Diffraction (10\,\%)
  and tissue property uncertainty (10\,\%) dominate; all approximation errors
  are below 3\,\%.  (b)~Worst case: 6\,GHz, finger illumination, single
  line-of-sight path.  Near-field proximity (2\,cm) dominates; the large
  diffraction error reflects the sub-wavelength size parameter.}
\label{fig:error-budget}
\end{figure}

% ============================================================
\section{Computational framework}\label{sec:computation}
% ============================================================

\subsection{Matrix form and neural-network interpretation}

Discretise the body surface into $M_{\mathrm{tri}}$ triangles with
normals assembled as rows of $\mathbf{N}\in\mathbb{R}^{M_{\mathrm{tri}}\times 3}$,
and collect the $N$ incident-wave directions as columns of
$\mathbf{K}\in\mathbb{R}^{3\times N}$.  Define
$\mathbf{M}=\mathbf{N}\mathbf{K}\in\mathbb{R}^{M_{\mathrm{tri}}\times N}$:
entry $\mu_{ji}=\nhat_j\cdot(-\khat_i)$ is the incidence cosine of path
$i$ at triangle $j$.  Let $\mathbf{s}=(S_1,\dots,S_N)^\top$ be the
path power densities.  The absorbed power density map is then
\begin{equation}\label{eq:Sab-mat}
  \mathbf{S}_{\mathrm{ab}} = T_0\;\ReLU(\mathbf{M})\;\mathbf{s},
\end{equation}
a single matrix-vector multiply followed by a ReLU and a weighted sum.
This is exactly the forward pass of a single-hidden-layer neural network
with linear input layer $\mathbf{K}^\top$, ReLU activation, and output
weights $T_0\,\mathbf{s}$.  The analogy is structural, not approximate.

\subsection{Hierarchical computation}

\begin{table}[!htbp]
\centering
\caption{Three levels of the computational hierarchy.  All levels are
  exact within their scope.}\label{tab:hierarchy}
\begin{tabular}{llll}
\toprule
Level & Output & Formula & Cost \\
\midrule
Bound    & Upper bound on $P_{\mathrm{abs}}$ & $\Sinc\,T_0\,A/2$ & $O(1)$ \\
Aggregate & Exact $P_{\mathrm{abs}}$, $\SARwb$ & $T_0\,\frac{\Aab}{4}\sum_i S_i D(\khat_i)$ & $O(N)$ \\
Spatial  & Full $\Sab(\rr)$ map & $T_0\,\ReLU(\mathbf{M})\,\mathbf{s}$ & $O(M_{\mathrm{tri}} N)$ \\
\bottomrule
\end{tabular}
\end{table}

The aggregate level suffices for whole-body SAR compliance checks.
For $M_{\mathrm{tri}}\sim 10^4$ triangles and $N\sim 10^2$ paths, it
is roughly $10^4$ times faster than the spatial computation.

\subsection{Precomputed body response}

All geometric quantities ($\eta(\rr)$, $\Aab$, $D(\khat)$) are
computed once from the body mesh, independent of the electromagnetic
environment.  At run time, the dosimetry computation is a table lookup
and a dot product.  The exposure fraction $\eta(\rr)$ is the ambient
occlusion map, which GPU rasterisers compute at interactive frame rates.

\subsection{Diffuse environment limit}

In a richly scattered indoor environment where the angular power
spectrum is approximately isotropic with total power density
$S_{\mathrm{total}}$, the directivity~$D$ averages to unity and
\cref{eq:gen-cauchy} gives directly:
\begin{equation}\label{eq:diffuse}
  P_{\mathrm{abs}} = \frac{S_{\mathrm{total}}\,T_0\,\Aab}{4},
  \qquad
  \Sab(\rr) = \frac{S_{\mathrm{total}}\,T_0}{4}\,\eta(\rr).
\end{equation}
The precomputed exposure fraction \emph{is} the dosimetry map.  No
per-direction computation is required.

\subsection{Near-field and point-source extension}

For a source at distance $d>3\lambda$ (locally plane-wave regime),
the geometric absorption law holds pointwise with spatially varying
inputs:
\begin{equation}\label{eq:Sab-point}
  \Sab(\rr) = \frac{P_t\,G\!\bigl(\khat(\rr)\bigr)}{4\pi\,d(\rr)^2}
    \cdot T_0 \cdot
    \ReLU\!\bigl[\nhat(\rr)\cdot\bigl(-\khat(\rr)\bigr)\bigr],
\end{equation}
where $P_t$ is the total radiated power, $G(\khat)$ is the antenna gain,
and $d(\rr)=|\rr-\rr_s|$ is the source-to-surface distance.  Integrating
over the illuminated surface gives the \emph{solid-angle formula}:
\begin{equation}\label{eq:solid-angle}
  \boxed{P_{\mathrm{abs}}
    = \frac{P_t\,T_0}{4\pi}\;\Omega_{\mathrm{body}}(\rr_s),}
\end{equation}
where $\Omega_{\mathrm{body}}(\rr_s)$ is the solid angle subtended by
the visible, front-facing body surface as seen from the source.  This
is the emissive radiative heat-transfer view-factor formula, with
$T_0$ replacing the emissivity.

% ============================================================
\partheading{Part~II: Generalisations}
{Sections~\ref{sec:polarisation}--\ref{sec:sub6}: extensions beyond
Part~I towards arbitrary polarisation and the sub-6\,GHz band.}
% ============================================================

% ============================================================
\section{Polarisation-aware dosimetry}\label{sec:polarisation}
% ============================================================

Part~I set the polarisation correction to zero.  The full correction is
handled via the Stokes representation.

\subsection{Absorption Stokes vector}

The polarisation correction in \cref{eq:Teff-decomp} depends on $q(\rr)$,
which in turn depends on the orientation of the incident polarisation
ellipse relative to the local TE/TM basis at each surface point.
Integrating \cref{eq:Sab-exact} over the body surface using the exact
$\Teff$ from \cref{eq:Teff-decomp} yields:
\begin{equation}\label{eq:Pabs-Stokes}
  \boxed{P_{\mathrm{abs}} = \mathbf{m}(\khat)\cdot\mathbf{s}_{\mathrm{inc}},}
\end{equation}
where $\mathbf{s}_{\mathrm{inc}}=(S_0,S_1,S_2,S_3)^\top$ is the Stokes
vector of the incident wave and the \emph{absorption Stokes vector} of the
body is
\begin{equation}\label{eq:m-def}
  \mathbf{m}(\khat) = \bigl(\tilde{A}_\perp,\;
    -\tfrac{1}{2}B_c,\;
    -\tfrac{1}{2}B_s,\;
    0\bigr)^\top.
\end{equation}
Here $\tilde{A}_\perp(\khat)\equiv\int_{\Sigma_+}\Tavg(\theta)\cos\theta\,\diff A$
is the \emph{Fresnel-weighted projected area} (approximately $T_0\Aperp$
but exact), and $\tilde{B}(\khat)=B_c+iB_s\equiv\int_{\Sigma_+}\Delta
T(\theta)\cos\theta\,e^{2i\alpha(\rr)}\,\diff A$ is the \emph{polarisation
response function}, where $\alpha(\rr)$ is the angle between the local
TE direction and the global reference polarisation basis.  Both
$\tilde{A}_\perp$ and $\tilde{B}$ are precomputable from the body mesh
and Fresnel coefficients, and are independent of the incident wave's
polarisation state.

Three properties follow immediately.  First, $S_3=0$ in the fourth
component: circular polarisation handedness does not affect total
absorbed power.  Second, for unpolarised light ($\mathbf{s}=S_0(1,0,0,0)^\top$),
$P_{\mathrm{abs}}=S_0\tilde{A}_\perp\approx S_0\,T_0\Aperp$,
recovering Part~I.  Third, for multiple sources the computation is
$O(N)$ via $P_{\mathrm{abs}}=\sum_i\mathbf{m}(\khat_i)\cdot\mathbf{s}_i$.

\subsection{Cylinder bound}

Define the \emph{polarisation directivity}
$D_B(\khat)\equiv|\tilde{B}(\khat)|/(2\tilde{A}_\perp(\khat))$.  For
linearly polarised illumination at the worst orientation,
$P_{\mathrm{abs}}=S_0\tilde{A}_\perp(1\pm D_B)$: the fractional
correction magnitude is~$D_B$.

On a sphere $D_B=0$ exactly, because the $e^{2i\alpha}$ phase rotates
through the full circle$[0,2\pi)$ uniformly over the illuminated
hemisphere and its integral vanishes.  The geometric extreme in the
opposite direction is an infinite vertical cylinder illuminated at
broadside, where $\alpha$ is the same at every surface point.  This
gives the \emph{cylinder bound}:
\begin{equation}
  D_B^{\mathrm{cyl}}
  = \frac{\int_0^1\Delta T(\mu)\,\mu\,\diff\mu}%
         {2\int_0^1\Tavg(\mu)\,\mu\,\diff\mu}.
\end{equation}
For skin at 28\,GHz, $D_B^{\mathrm{cyl}}=27.9\,\%$, an upper bound for
any convex body and any direction.  On the Thelonious human phantom,
the maximum over all directions is $D_B=16.0\,\%$, occurring near
lateral illumination (where cylindrical limbs dominate).

\subsection{Multipath convergence}

In an indoor environment with $N$ paths whose polarisation orientations
are mutually independent and uniformly distributed, the standard deviation
of the fractional polarisation error decreases as
\begin{equation}
  \frac{\sigma[\Delta P_{\mathrm{pol}}]}{P_{\mathrm{unpol}}}
  = \frac{D_B(\khat)}{\sqrt{2N}}.
\end{equation}
With $D_B\le 16\,\%$ and $N\ge 20$ paths the error is below 2.5\,\%,
comparable to the pseudo-Brewster approximation.  In a reverberation
chamber ($N\to\infty$) the correction vanishes exactly.  The
\emph{unpolarised geometric absorption law is conservative for
vertically polarised base-station illumination}: because horizontal
limbs are purely TE, and $T_s(\theta)\le T_0$ for all $\theta$, the
simplified formula overestimates limb absorption.

% ============================================================
\section{Extension below 6\,GHz}\label{sec:sub6}
% ============================================================

\subsection{Exact flux-averaged transmission}

The pseudo-Brewster approximation $\Tavg(\theta)\approx T_0$ weakens at
sub-10\,GHz frequencies.  However, the direction-averaged total power
formula is exact if we replace $T_0$ by the \emph{flux-averaged
transmission}:
\begin{equation}\label{eq:Tbar}
  \Tbar(f) \equiv 2\int_0^1 \Tavg(\mu,f)\;\mu\;\diff\mu,
\end{equation}
which is the fraction of incident power absorbed by a convex body in an
isotropic, unpolarised field.  It is computed once per frequency from
the Fresnel coefficients and depends only on $\ntilde(f)$.  With this
substitution, \cref{eq:gen-cauchy} becomes:
\begin{equation}\label{eq:Cauchy-exact}
  \boxed{\langle P_{\mathrm{abs}}\rangle = \Sinc\,\Tbar\,\frac{\Aab}{4},}
\end{equation}
exact at any frequency above ${\sim}100$\,MHz (where the body is
electromagnetically opaque).

\begin{table}[!htbp]
\centering
\caption{Normal-incidence transmission $T_0$, exact flux-averaged
  transmission $\Tbar$, and ratio $R=T_0/\Tbar$ for skin
  (IT'IS v5.0~\cite{ITISv5}).  The ratio $R$ crosses unity at
  40\,GHz, the pseudo-Brewster sweet spot.}\label{tab:Tbar}
\begin{tabular}{rcccr}
\toprule
$f$\,(GHz) & $|\ntilde|$ & $T_0$ & $\Tbar$ & $R$ \\
\midrule
  0.9 & 6.70 & 0.445 & 0.465 & 0.957 \\
  6.0 & 6.07 & 0.481 & 0.497 & 0.968 \\
 28   & 4.84 & 0.536 & 0.543 & 0.988 \\
 40   & 4.30 & 0.570 & 0.571 & 1.000 \\
 60   & 3.68 & 0.622 & 0.613 & 1.015 \\
100   & 3.01 & 0.701 & 0.677 & 1.035 \\
\bottomrule
\end{tabular}
\end{table}

Below 40\,GHz the simplified formula (using $T_0$) underestimates
absorbed power and is therefore conservative for compliance.  Above
40\,GHz the overestimate is at most 3.5\,\%.

\subsection{Validity summary}

The local surface map $\Sab(\rr)\approx \Sinc T_0\ReLU(\mu)$ loses
physical meaning below 6\,GHz, when the skin depth approaches the surface
layer thickness and volumetric absorption and multi-layer resonances
(Fabry--P\'{e}rot effects between skin, fat, and muscle) become significant.
Total power and compliance results remain valid at all frequencies above
100\,MHz via $\Tbar$.

% ============================================================
\partheading{Part~III: Coherent MIMO Illumination}
{Sections~\ref{sec:coherent-setup}--\ref{sec:ecbf}: the framework
extended to beamformed sources with controlled complex weights.}
% ============================================================

Parts~I and~II assume incoherent illumination: paths arrive with
random relative phases so that cross-terms average to zero.  In coherent
MIMO beamforming, $M$ antenna elements radiate with controlled complex
weights and their fields add coherently.  For $N$ co-phased paths the
coherent peak can exceed the incoherent sum by a factor of $N$; the
Part~I--II framework cannot capture this.

% ============================================================
\section{Setup}\label{sec:coherent-setup}
% ============================================================

A base station has $M$ antenna elements with precoding vector
$\mathbf{x}=[x_1,\ldots,x_M]^\top\in\mathbb{C}^M$,
$\|\mathbf{x}\|^2=P$ (total radiated power).  Between the base station
and a user body, the propagation environment creates $N$ paths.  Path
$n$ is characterised by:
\begin{itemize}[noitemsep]
  \item its direction of arrival at the body $\hat{\mathbf{k}}_n$;
  \item its originating antenna element $j(n)\in\{1,\ldots,M\}$;
  \item its \emph{polarisation-amplitude vector}
        $\boldsymbol{\psi}_n\in\mathbb{C}^3$
        ($\boldsymbol{\psi}_n\perp\hat{\mathbf{k}}_n$,
        units: $\mathrm{V\,m^{-1}\,W^{-1/2}}$), which encodes the path's
        complex amplitude, polarisation state, and all propagation interactions
        (reflections, scattering, distance attenuation).  It is obtained
        directly from a ray tracer.
\end{itemize}
The \emph{field channel matrix} at surface point $\mathbf{r}$ is
\begin{equation}\label{eq:G-def}
  \mathbf{G}(\mathbf{r})
  = \bigl[\mathbf{g}_1(\mathbf{r}),\ldots,\mathbf{g}_M(\mathbf{r})\bigr]
  \in\mathbb{C}^{3\times M},
  \qquad
  \mathbf{g}_j(\mathbf{r}) = \sum_{n:\,j(n)=j}
    \boldsymbol{\psi}_n\,e^{-ik_0\hat{\mathbf{k}}_n\cdot\mathbf{r}},
\end{equation}
so that the total electric field is $\mathbf{E}(\mathbf{r})=\mathbf{G}(\mathbf{r})\mathbf{x}$.

% ============================================================
\section{Coherent absorption law}\label{sec:coherent-law}
% ============================================================

\subsection{Fresnel transmission in the coherent regime}

The pseudo-Brewster approximation ($\Teff\approx T_0$) used in Part~I
replaced angle-dependent complex Fresnel amplitude coefficients
$t_s(\theta)$, $t_p(\theta)$ with the real constant $\sqrt{T_0}$.  In
the coherent regime this fails: absorbed power involves cross-products
$t_{s,n}^*t_{s,n'}$ between paths at different incidence angles, and
these carry angle-dependent phases that the constant approximation
destroys.  Exact Fresnel amplitudes are therefore retained throughout
Part~III.

The \emph{Fresnel transmission operator} for path $n$ at surface
point~$\mathbf{r}$ is
\begin{equation}\label{eq:Fn}
  \mathbf{F}_n(\mathbf{r}) =
  \Bigl(
    t_{s,n}\;\hat{e}_{s,n}\hat{e}_{s,n}^\top
    + t_{p,n}\;\hat{e}_{p,n}\hat{e}_{p,n}^\top
  \Bigr)H\!\bigl(\mu_n(\mathbf{r})\bigr),
\end{equation}
where $\hat{e}_{s,n}$ and $\hat{e}_{p,n}$ are the local TE and TM unit
vectors for path $n$ at $\mathbf{r}$, $t_{s,n}$ and $t_{p,n}$ are the
corresponding Fresnel amplitude transmission coefficients (scalars,
complex), and $\mu_n(\mathbf{r})=\nhat(\mathbf{r})\cdot(-\hat{\mathbf{k}}_n)$
is the incidence cosine.  The operator $\mathbf{F}_n$ is dimensionless
and rank-2: it projects $\boldsymbol{\psi}_n$ onto the TE/TM basis and
applies the correct transmission amplitude to each component.  The
$H(\mu_n)$ factor zeros back-facing paths ($\mu_n\le 0$), the exact
analogue of the ReLU visibility gate in Part~I.

\subsection{Depth integral and two approximations}

The absorbed power density at $\mathbf{r}$ is the Ohmic loss integrated
over depth $z$ into the tissue:
\begin{equation}\label{eq:Sab-depth}
  \Sab(\mathbf{r}) = \frac{\sigma}{2}\int_0^\infty
  \Bigl|\sum_n \mathbf{F}_n(\mathbf{r})\boldsymbol{\psi}_n
    x_{j(n)} e^{-ik_0\hat{\mathbf{k}}_n\cdot\mathbf{r}}
    e^{-ik_0\xi_n z}\Bigr|^2\diff z,
\end{equation}
where $\sigma$ is the tissue conductivity and $\xi_n$ is the (complex)
normal wave-vector in the tissue for path $n$.  Expanding the squared
modulus produces cross-terms $(n,n')$, each weighted by a
\emph{depth coupling factor} $\Lambda_{nn'}=\int_0^\infty
e^{-(\alpha_n+\alpha_{n'})z}e^{i(\beta_n-\beta_{n'})z}\diff z$, where
$k_0\xi_n=\beta_n-i\alpha_n$ defines the per-path decay rate $\alpha_n$
and phase rate $\beta_n$ inside the tissue.

Two approximations are applied.

\begin{proposition}[Approximation~1: incident TM direction]
The physically exact transmitted TM field points along the refracted
direction inside the tissue, not the incident direction.  Replacing
the refracted by the incident direction introduces an error
$O(1/|\ntilde|^2)\lesssim 4\,\%$ on cross-terms at 28\,GHz.  Self-terms
are unaffected.
\end{proposition}

\begin{proposition}[Approximation~2: universal depth coupling]
The normalised depth coupling
$\Gamma_{nn'}\equiv\Lambda_{nn'}/\sqrt{\Lambda_{nn}\Lambda_{n'n'}}$
satisfies $|1-\Gamma_{nn'}|\le 0.44\,\%$ for all incidence-angle pairs
at 28\,GHz.  This near-universality follows from the high refractive
index $|\ntilde|\approx 4.8$, which compresses all refracted angles to
within ${\sim}10^\circ$ of normal, making $\alpha_n$ and $\beta_n$
nearly angle-independent.  We set $\Gamma_{nn'}=1$ for all pairs.
\end{proposition}

Under these two approximations, the depth coupling factors factorise:
$\Lambda_{nn'}\approx\sqrt{\Lambda_{nn}\Lambda_{n'n'}}$, and the double
sum \cref{eq:Sab-depth} becomes a perfect square.

\subsection{Main theorem}

Define the \emph{exposure channel matrix}:
\begin{equation}\label{eq:Gtilde-def}
  \tilde{\mathbf{G}}(\mathbf{r})
  = \bigl[\tilde{\mathbf{g}}_1(\mathbf{r}),\ldots,\tilde{\mathbf{g}}_M(\mathbf{r})\bigr]
  \in\mathbb{C}^{3\times M},
  \quad
  \tilde{\mathbf{g}}_j(\mathbf{r}) = \sum_{n:\,j(n)=j}
    \sqrt{\frac{\sigma}{4\alpha_n}}\;
    \mathbf{F}_n(\mathbf{r})\,\boldsymbol{\psi}_n\,
    e^{-ik_0\hat{\mathbf{k}}_n\cdot\mathbf{r}}.
\end{equation}
Each column of $\tilde{\mathbf{G}}$ is the effective field contribution
from one transmit element after Fresnel filtering, depth-decay weighting,
and phase propagation.  It is the absorption-equivalent of the field
channel matrix $\mathbf{G}$, but weighted to convert field amplitude
directly into absorbed power.

\begin{theorem}[Coherent absorption law]\label{thm:coherent-law}
Under Approximations~1 and~2,
\begin{equation}\label{eq:coherent-law}
  \boxed{\Sab(\mathbf{r}) = \|\tilde{\mathbf{G}}(\mathbf{r})\,\mathbf{x}\|^2.}
\end{equation}
This holds for any precoding vector $\mathbf{x}$.
\end{theorem}

\begin{corollary}[Incoherent limit]
When paths have mutually random phases, cross-terms average to zero and
\cref{eq:coherent-law} reduces to the Part~I multi-source formula.
\end{corollary}

% ============================================================
\section{Exposure operator}\label{sec:Q}
% ============================================================

Integrating \cref{eq:coherent-law} over the body surface:
\begin{equation}\label{eq:Q-def}
  P_{\mathrm{abs}} = \mathbf{x}^H\,\mathbf{Q}\,\mathbf{x},
  \qquad
  \boxed{\mathbf{Q} \equiv \int_\Sigma
    \tilde{\mathbf{G}}(\mathbf{r})^H\,\tilde{\mathbf{G}}(\mathbf{r})\,\diff A
    \;\in\;\mathbb{C}^{M\times M}.}
\end{equation}
The \emph{exposure operator} $\mathbf{Q}$ is Hermitian
positive-semidefinite and provides a complete, finite-dimensional
description of the body's absorption response to any precoder.  Its
eigendecomposition $\mathbf{Q}=\sum_k\lambda_k\mathbf{v}_k\mathbf{v}_k^H$ reveals
the \emph{exposure modes}: $\mathbf{v}_1$ is the precoder that maximises
total absorption; $\lambda_1$ is the maximum absorbed power per unit
transmit power; $\mathrm{tr}(\mathbf{Q})$ is the expected absorption
under isotropic transmission.  This plays the same role in Part~III as
the absorption directivity $D(\khat)$ played in Part~I.

Under maximum-ratio transmission (MRT, $\mathbf{x}=\sqrt{P}\,\mathbf{h}^*/\|\mathbf{h}\|$
for UE channel vector $\mathbf{h}$), absorbed power is
$P_{\mathrm{abs}}^{\mathrm{MRT}}=P\,\rho\,\lambda_{\max}(\mathbf{Q})$, where
the \emph{exposure-signal alignment}
\begin{equation}\label{eq:rho-def}
  \boxed{\rho \equiv \frac{\mathbf{h}^\top\mathbf{Q}\,\mathbf{h}^*}
    {\|\mathbf{h}\|^2\,\lambda_{\max}(\mathbf{Q})} \;\in\;[0,1]}
\end{equation}
is the single number quantifying the exposure cost of MRT.  It is the
normalised Rayleigh quotient of $\mathbf{Q}$ in the direction
$\mathbf{h}^*$.  A small $\rho$ means MRT naturally avoids illuminating
the body's primary absorption mode; body geometry provides inherent
protection.

The signal branch sums scalars ($y\propto\sum_j|h_j|^2$); the
exposure branch at $\mathbf{r}$ sums vectors
$h_j^*\,\tilde{\mathbf{g}}_j(\mathbf{r})$.  Three structural
differences prevent the body hotspot from being a copy of the signal
focus.  First, transmitted polarisations of co-originating paths
differ at $\mathbf{r}$ and can partially cancel even when all scalar
magnitudes $|h_j^*||\tilde{g}_j|$ are equal.  Second, the UE antenna
projects each path to a scalar via $\mathbf{C}_R^H\boldsymbol{\psi}_n$,
while the exposure channel retains the full vector structure via
$\mathbf{F}_n(\mathbf{r})\boldsymbol{\psi}_n$; a path at grazing
incidence contributes strongly to the received signal but weakly to
body absorption.  Third, moving from the UE position
$\mathbf{r}_\mathrm{UE}$ to a body-surface point $\mathbf{r}$,
each path acquires the additional phase
$e^{-ik_0\hat{\mathbf{k}}_n\cdot(\mathbf{r}-\mathbf{r}_\mathrm{UE})}$,
which degrades the coherent alignment built up at the UE.

% ============================================================
\section{Exposure-constrained beamforming}\label{sec:ecbf}
% ============================================================

The standard communications problem — maximise received signal power
subject to a body absorption limit — becomes:
\begin{equation}\label{eq:QCQP}
  \max_{\mathbf{x}}\;|{\mathbf{h}^\top\mathbf{x}}|^2
  \quad\text{s.t.}\quad
  \mathbf{x}^H\mathbf{Q}\,\mathbf{x}\le P_{\mathrm{abs}}^{\max},
  \quad\|\mathbf{x}\|^2\le P.
\end{equation}
Both constraints are convex (PSD matrices) with a non-convex quadratic
objective.  For QCQPs with at most two quadratic constraints, the
KKT point is a global optimum (S-procedure).  The optimal precoder is
\begin{equation}\label{eq:optimal-x}
  \boxed{\mathbf{x}^\star = \sqrt{P}\;
    \frac{(\lambda\mathbf{Q}+\nu\mathbf{I})^{-1}\mathbf{h}^*}
         {\|(\lambda\mathbf{Q}+\nu\mathbf{I})^{-1}\mathbf{h}^*\|},}
\end{equation}
with Lagrange multipliers $\lambda\ge 0$ (exposure constraint) and
$\nu\ge 0$ (power constraint) found by complementary slackness.  The
matrix $(\lambda\mathbf{Q}+\nu\mathbf{I})^{-1}$ regularises MRT in the
$\mathbf{Q}$-eigenspace: directions with large eigenvalue (high
absorption modes) are suppressed by $\sim 1/(\lambda\lambda_k)$.
Null-space directions of $\mathbf{Q}$ (zero body absorption) are
unconstrained.

For MU-MIMO with $K$ simultaneous users, each user $u$ has an exposure
operator $\mathbf{Q}^{(u)}$ built from the paths arriving at their body.
The expected absorbed power in person $u$'s body from all $K$ data
streams sums incoherently (data symbols of different users are
uncorrelated):
\begin{equation}\label{eq:Pabs-MU}
  P_{\mathrm{abs}}^{(u)} = \sum_{k=1}^K \mathbf{w}_k^H\,\mathbf{Q}^{(u)}\,\mathbf{w}_k,
\end{equation}
where $\mathbf{w}_k$ is the precoder for user~$k$.  The per-person
exposure constraints are quadratic with PSD matrices and fit directly
into WMMSE or successive-convex-approximation solvers.

% ============================================================
\section{Conclusion}\label{sec:conclusion}
% ============================================================

The absorbed power density on biological tissue from 100\,MHz to
100\,GHz admits a hierarchy of closed-form expressions.  Two physical
facts underpin the framework.

First, pseudo-Brewster compensation makes the unpolarised Fresnel
transmission approximately angle-independent over 50\,\% of the
electromagnetic spectrum relevant to wireless communication
(10--60\,GHz), collapsing the material factor to a scalar~$T_0$.
Replacing $T_0$ with the exact flux-average~$\Tbar$ makes
direction-averaged results valid at any frequency above 100\,MHz.

Second, the small skin depth above 6\,GHz confines absorption to the
surface, converting the full volumetric problem into a geometric one.
Total absorbed power is $\Sinc T_0$ times the projected area.  The
direction-average satisfies Cauchy's formula $\langle P\rangle =
\Sinc\Tbar\Aab/4$, with the absorption area $\Aab$ identical to
the ambient-occlusion integral of computer graphics.  For multiple
sources the formula is a single-hidden-layer ReLU network, directly
differentiable through any ray-tracing engine.

Polarisation corrections are bounded universally by 27.9\,\% (cylinder
limit) and suppressed below 2.5\,\% by multipath diversity for
$N\ge 20$ paths.

For coherent MIMO beamforming, two small approximations (each
$\lesssim 5\,\%$) collapse the coherent depth integral to the squared
norm $\Sab(\mathbf{r})=\|\tilde{\mathbf{G}}(\mathbf{r})\mathbf{x}\|^2$.
Integrating over the body yields the exposure operator $\mathbf{Q}$,
whose eigendecomposition reveals the body's absorption modes and enables
a closed-form exposure-constrained precoder.  The total computational
cost is $O(M_{\mathrm{tri}}\,M^2)$ operations, reducing a
$10^{12}$-cell volumetric simulation to milliseconds on a GPU.


\begingroup
\raggedright
\begin{thebibliography}{99}

% ---- Standards and guidelines ----
\bibitem{ICNIRP2020}
  ICNIRP, ``Guidelines for limiting exposure to electromagnetic fields
  (100\,kHz to 300\,GHz),'' \textit{Health Phys.}, vol.~118, no.~5,
  pp.~483--524, 2020.

% ---- Tissue dielectric properties ----
\bibitem{ITISv5}
  IT'IS Foundation, ``Tissue properties database v5.0,'' 2024.
  [Online]. Available: \url{https://itis.swiss/virtual-population/tissue-properties/database/}

\bibitem{Gabriel1996}
  S.~Gabriel, R.~W.~Lau, and C.~Gabriel, ``The dielectric properties of
  biological tissues: III.\ Parametric models for the dielectric
  spectrum of tissues,'' \textit{Phys.\ Med.\ Biol.}, vol.~41, no.~11,
  pp.~2271--2293, 1996.

\bibitem{AlekseevZiskin2007}
  S.~I.~Alekseev and M.~C.~Ziskin, ``Human skin permittivity determined
  by millimeter wave reflection measurements,''
  \textit{Bioelectromagnetics}, vol.~28, no.~5, pp.~331--339, 2007.

% ---- EM dosimetry ----
\bibitem{Christ2006}
  A.~Christ, T.~Samaras, A.~Klingenb\"{o}ck, and N.~Kuster,
  ``Characterization of the electromagnetic near-field absorption in
  layered biological tissue in the frequency range from 30\,MHz to
  300\,GHz,'' \textit{Phys.\ Med.\ Biol.}, vol.~51, no.~19,
  pp.~4951--4965, 2006.

\bibitem{Dimbylow2002}
  P.~J.~Dimbylow, ``Fine resolution calculations of SAR in the human
  body for frequencies up to 3\,GHz,'' \textit{Phys.\ Med.\ Biol.},
  vol.~47, no.~16, pp.~2835--2846, 2002.

\bibitem{Colombi2018}
  D.~Colombi, B.~Thors, C.~Torne\-vik, and Q.~Balzano, ``RF energy
  absorption by biological tissues in close proximity to millimeter-wave
  5G wireless equipment,'' \textit{IEEE Access}, vol.~6, pp.~4974--4981,
  2018.

% ---- Pseudo-Brewster / Fresnel physics ----
\bibitem{Azzam2015}
  R.~M.~A.~Azzam, ``High-index dielectric substrates with nearly
  constant reflectance for incident unpolarized or circularly polarized
  light over a wide range of incidence angles,'' \textit{J.\ Mod.\
  Opt.}, vol.~62, no.~18, pp.~1503--1507, 2015.

\bibitem{BornWolf1999}
  M.~Born and E.~Wolf, \textit{Principles of Optics}, 7th (expanded)
  ed.\hskip1em Cambridge Univ.\ Press, 1999.

% ---- Room electromagnetics and diffuse-field dosimetry ----
\bibitem{Bamba2014}
  A.~Bamba, W.~Joseph, G.~Vermeeren, E.~Tanghe, D.~P.~Gaillot,
  J.~B.~Andersen, J.~{\O}.~Nielsen, M.~Lienard, and L.~Martens,
  ``Validation of experimental whole-body SAR assessment method in a
  complex indoor environment,'' \textit{Bioelectromagnetics}, vol.~34,
  no.~2, pp.~122--132, 2013.

% ---- Integral geometry ----
\bibitem{Cauchy1841}
  A.-L.~Cauchy, ``Note sur divers th\'{e}or\`{e}mes relatifs \`{a} la
  rectification des courbes et \`{a} la quadrature des surfaces,''
  \textit{Comptes Rendus Acad.\ Sci.}, vol.~13, pp.~1060--1065, 1841.

% ---- Ambient occlusion / computer graphics ----
\bibitem{Zhukov1998}
  S.~Zhukov, A.~Iones, and G.~Kronin, ``An ambient light illumination
  model,'' in \textit{Rendering Techniques~'98} (Proc.\ Eurographics
  Workshop on Rendering), G.~Drettakis and N.~Max, Eds.\hskip1em
  Springer, 1998, pp.~45--55.

% ---- Neural networks and ReLU ----
\bibitem{NairHinton2010}
  V.~Nair and G.~E.~Hinton, ``Rectified linear units improve restricted
  Boltzmann machines,'' in \textit{Proc.\ 27th Int.\ Conf.\ Machine
  Learning (ICML)}, 2010, pp.~807--814.

% ---- Body surface area / anthropometry ----
\bibitem{DuBois1916}
  D.~Du~Bois and E.~F.~Du~Bois, ``Clinical calorimetry: Tenth paper.
  A formula to estimate the approximate surface area if height and
  weight be known,'' \textit{Arch.\ Intern.\ Med.}, vol.~17,
  pp.~863--871, 1916.

% ---- Mie theory ----
\bibitem{BohrenHuffman1983}
  C.~F.~Bohren and D.~R.~Huffman, \textit{Absorption and Scattering of
  Light by Small Particles}.\hskip1em Wiley, 1983.

% ---- Ray tracing ----
\bibitem{SionnaRT}
  J.~Hoydis \textit{et al.}, ``Sionna: An open-source library for
  next-generation physical layer research,'' 2023. [Online]. Available:
  \url{https://nvlabs.github.io/sionna/}

\end{thebibliography}
\endgroup

\end{document}
